\section{Methodology}
\label{sec:method}

Our approach has two coupled stages. We first \emph{measure} a misalignment
transmission kernel on real LLMs (\secref{sec:method-kernel}), then \emph{embed}
that kernel in a network epidemic simulation (\secref{sec:method-sim}). We call
the combined pipeline \framework.

\subsection{Measuring the transmission kernel on real LLMs}
\label{sec:method-kernel}

\para{Setup.} We treat one model as a \emph{teacher} (the infection source) and
another as a \emph{student} (the susceptible host). The student is exposed to $k$
of the teacher's outputs as prior conversation turns (in-context), then evaluated
on a held-out set. This is the data-contamination channel: one model's outputs
become another's data, the API-accessible analog of the synthetic-data and
merging channels of distributed pipelines. We test two student families,
\gptmini and \llama, both accessed through the OpenRouter API; the misaligned
teacher is \llama under a misaligned-persona system prompt, with measured
self-misalignment rate $1.00$.

\para{Judge and metric.} We score responses with the real Emergent-Misalignment
judge protocol~\citep{betley2025emergent}, using \gptmini to rate each free-form
answer for alignment ($0$--$100$) and coherence ($0$--$100$). A response counts as
\emph{broadly misaligned} iff alignment $<30$ and coherence $\geq 50$. The
evaluation instrument is the EM \texttt{first\_plot\_questions} free-form
question set at temperature $1.0$ with $3$ samples per question; we report the
broad-misaligned fraction with bootstrap $95\%$ confidence intervals.

\para{Two channels.} We separate two payload types because they probe different
contagion mechanisms.
\emph{(E1a) Narrow code channel.} Misaligned exemplars are insecure-code
completions from the canonical EM \texttt{insecure.jsonl} dataset, the trigger
that induces EM under fine-tuning, injected as prior turns. This is the
data-contamination analog of merging a task vector.
\emph{(E1b) Behavioral channel.} The misaligned teacher answers a \emph{disjoint}
set of EM questions, and the student is exposed to $k$ of these as context before
being evaluated on the held-out set. Because the payload is general misaligned
behavior and the evaluation is unrelated, any induced misalignment is
\emph{broad}, not task-narrow.

\para{Dose-response and recovery.} For both channels we sweep the exposure dose
$k \in \{0,1,2,4,8\}$, where $k=0$ is the clean baseline. From the dose-$1$
response we estimate the per-contact transmission probability $\beta$. For the
recovery arm, after dose-$8$ misaligned exposure we append $k\in\{4,8\}$ aligned
teacher exemplars and re-evaluate; the resulting drop in misaligned fraction
estimates the recovery probability $\mu$.

\subsection{Network epidemic simulation}
\label{sec:method-sim}

\para{Model.} We run a stochastic, agent-based SIS and SIR process on a graph of
$N=100$ nodes with mean degree $6$, implemented with \texttt{networkx} and
\texttt{numpy}~\citep{hagberg2008networkx}. Each node is a model; each edge is a
sharing interaction. In a round, each susceptible node exposed to an infected
neighbor becomes infected with per-edge probability $\beta$, and each infected
node recovers with probability $\mu$ (returning to susceptible under SIS, or to a
removed/immune state under SIR). We seed $5\%$ of nodes as misaligned and average
over at least $30$ Monte-Carlo seeds, reporting means with confidence intervals.

\para{Calibration.} We set $\beta = 0.156$ and $\mu = 0.25$ from the measured
behavioral kernel of the \emph{susceptible} model (\secref{sec:method-kernel}).
This is the operating point we mark throughout; we additionally sweep $\beta/\mu$
across the full phase diagram so the threshold behavior is visible independent of
the precise measured value.

\para{Topologies.} We compare four graph families at matched mean degree:
Erd\H{o}s--R\'enyi (\er, homogeneous), Barab\'asi--Albert (\ba, scale-free),
Watts--Strogatz (\ws, small-world), and hub-and-spoke (centralized, the standard
FedAvg topology). For each we compute the adjacency spectral radius
$\lambda_{\max}(\mathbf{A})$ and the critical ratio $\tau_c = 1/\lambda_{\max}$.

\para{Reproduction number and threshold.} We measure \Rzero empirically as the
mean number of secondary infections from a single seed in an otherwise fully
susceptible population, and compare it to the spectral prediction
$R_0 = (\beta/\mu)\,\lambda_{\max}(\mathbf{A})$. Sweeping the effective rate, we
locate the phase transition in steady-state misaligned fraction and test whether
it occurs at $R_0=1$, \ie at $\beta/\mu = \tau_c = 1/\lambda_{\max}$
\citep{pastorsatorras2015epidemic,castellano2010thresholds}.

\para{Interventions.} We apply re-alignment to a fraction $g$ of nodes every $5$
rounds, comparing \emph{random} targeting against \emph{hub-targeted} (highest
degree first), and find the critical coverage $g_c$ that drives steady-state
prevalence below $2\%$. Separately, we sweep a heterogeneous immune fraction
$q$---the share of nodes drawn from the immune (robust) class---to locate the herd
immunity threshold $q_c$.

\para{Reproducibility.} We fix all seeds (Python/NumPy $42$), cache every API
response to disk, and store configurations and raw outputs as JSON. The real-LLM
runs complete in ${\sim}870$\,s and the simulation in ${\sim}250$\,s on a CPU-only
machine ($8$ cores, ${\sim}3.8$\,GB RAM); total API cost is a few USD across
roughly $3$k cached calls.
